\section{交织函子的逆}\label{sec:03.3.4}

对任意 $w\in W$ 与 $\lambda\in\h^*$，函子 $LI_w:D(\D_\lambda)\to D(\D_{w\lambda})$ 是范畴等价。它诱导有界凝聚复形所成子范畴之间的等价
$$LI_w:D^b_{coh}(\D_\lambda)\longrightarrow D^b_{coh}(\D_{w\lambda}).$$
本节要证明这些函子的拟逆的一个公式。

设 $\lambda\in\h^*$，$\theta=W\cdot\lambda$。记 $-\theta$ 为 $-\lambda$ 的 Weyl 群轨道。
\begin{annote}
底本此处作「记 $-\theta$ 为 $-\lambda$ 的 Weyl 群」，但 $-\theta$ 是轨道而不是群。本节稍后正要把 $U_\theta$ 的反环与 $U_{-\theta}$ 等同，可见 $-\theta$ 指的是 $-\lambda$ 所在的 Weyl 群轨道。[来源：底本 p.59，与下文 $U_\theta^\circ\cong U_{-\theta}$ 的用法对照]
\end{annote}

先回忆有界凝聚复形上的对偶函子 $\mathbb{D}:D^b_{coh}(\D_\lambda)\to D^b_{coh}(\D_{-\lambda})$。对任意复形 $V^\bullet$ 有
$$\mathbb{D}(V^\bullet)=R\Hom_{\D_\lambda}\bigl(V^\bullet,\mathbb{D}(\D_\lambda)\bigr)[\dim X].$$
这个对偶运算与张量积相容。

\begin{lemma}\label{lem:03.4.1}
对任意权 $\nu\in P(\Sigma)$，下述函子图表交换
\begin{center}
\begin{tikzcd}[column sep=4.6em,row sep=2.9em]
D^b_{coh}(\D_\lambda) \arrow[r,"\mathbb{D}"] \arrow[d,"(\nu)"'] & D^b_{coh}(\D_{-\lambda}) \arrow[d,"(-\nu)"]\\
D^b_{coh}(\D_{\lambda+\nu}) \arrow[r,"\mathbb{D}"'] & D^b_{coh}(\D_{-\lambda-\nu})
\end{tikzcd}
\end{center}
\end{lemma}

\begin{proof}
设 $V^\bullet$ 是凝聚 $\D_\lambda$-模的有界复形。则 $V^\bullet(\nu)$ 是凝聚 $\D_{\lambda+\nu}$-模的有界复形。此外，
{\small
\begin{align*}
\mathbb{D}(V^\bullet(\nu))
&=R\Hom_{\D_{\lambda+\nu}}\bigl(V^\bullet(\nu),\mathbb{D}(\D_{\lambda+\nu})\bigr)[\dim X]\\
&=R\Hom_{\D_{\lambda+\nu}}\bigl(\calO(\nu)\otimes_{\calO_X}V^\bullet,\ \calO(\nu)\otimes_{\calO_X}\mathbb{D}(\D_\lambda)\otimes_{\calO_X}\calO(-\nu)\bigr)[\dim X]\\
&=R\Hom_{\D_{\lambda+\nu}}\bigl(\calO(\nu)\otimes_{\calO_X}V^\bullet,\ \calO(\nu)\otimes_{\calO_X}\mathbb{D}(\D_\lambda)\bigr)[\dim X]\otimes_{\calO_X}\calO(-\nu)\\
&=R\Hom_{\D_\lambda}\bigl(V^\bullet,\mathbb{D}(\D_\lambda)\bigr)[\dim X]\otimes_{\calO_X}\calO(-\nu)
=\mathbb{D}(V^\bullet)(-\nu).
\end{align*}}
\end{proof}

如 3.4 那样，可见范畴 $\mathcal{M}_{coh}(\D_\lambda)$ 有足够多的投射对象。设 $P$ 是 $\mathcal{M}_{coh}(\D_\lambda)$ 中的投射对象。则由谱序列
$$H^p\bigl(X,\Ext^q_{\D_\lambda}(P,\D_\lambda)\bigr)\Longrightarrow\Ext^{p+q}_{\D_\lambda}(P,\D_\lambda)$$
推知，当 $p>0$ 时 $H^p\bigl(X,\Hom_{\D_\lambda}(P,\D_\lambda)\bigr)=0$；也就是说，$\Hom_{\D_\lambda}(P,\D_\lambda)$ 关于 $\Gamma(X,-)$ 零调。

考虑从 $D^-_{coh}(\D_\lambda)$ 映入 $D^+_{coh}(\D_{-\lambda})$ 的函子 $V^\bullet\mapsto R\Hom_{\D_\lambda}\bigl(V^\bullet,\mathbb{D}(\D_\lambda)\bigr)$，以及从 $D^+_{coh}(\D_{-\lambda})$ 映入 $D^+(U_\theta)$ 的函子 $R\Gamma$。上面的注记说明
$$R\Gamma\bigl(R\Hom_{\D_\lambda}(V^\bullet,\mathbb{D}(\D_\lambda))\bigr)=R\Hom_{\D_\lambda}\bigl(V^\bullet,\mathbb{D}(\D_\lambda)\bigr).$$
由此得到下述结果。

\begin{lemma}\label{lem:03.4.2}
从 $D^b_{coh}(\D_\lambda)$ 映入 $D^b(U_\theta)$ 的函子之间有同构
$$R\Gamma(\mathbb{D}(V^\bullet))=R\Hom_{\D_\lambda}\bigl(V^\bullet,\mathbb{D}(\D_\lambda)\bigr)[\dim X].$$
\end{lemma}

设 $\theta$ 是正则轨道。对这种 $\theta$，环 $U_\theta$ 的同调维数有限。此外，$U(\g)$ 的主反自同构诱导环同构 $U_\theta^\circ\cong U_{-\theta}$。记 $D^b_{fg}(U_\theta)$ 为有限生成 $U_\theta$-模的导出范畴。定义反变对偶函子
$$\mathbb{D}_{alg}(V^\bullet)=R\Hom_{U_\theta}\bigl(V^\bullet,\mathbb{D}(U_\theta)\bigr),$$
它从 $D^b_{fg}(U_\theta)$ 映入 $D^b_{fg}(U_{-\theta})$。显然 $\mathbb{D}_{alg}^2\cong 1$。

设 $V^\bullet$ 是上方有界的有限生成 $U_\theta$-模复形。则存在上方有界的复形 $F^\bullet$ 与复形态射 $F^\bullet\to V^\bullet$，其中 $F^\bullet$ 由有限秩自由 $U_\theta$-模组成。于是
{\small
\begin{align*}
R\Gamma\bigl(\mathbb{D}(L\Delta_\lambda(V^\bullet))\bigr)
&=R\Gamma\bigl(\mathbb{D}(\Delta_\lambda(F^\bullet))\bigr)
=R\Hom_{\D_\lambda}\bigl(\Delta_\lambda(F^\bullet),\mathbb{D}(\D_\lambda)\bigr)[\dim X]\\
&=\Hom^\bullet_{\D_\lambda}\bigl(\Delta_\lambda(F^\bullet),\mathbb{D}(\D_\lambda)\bigr)[\dim X]
=\Hom^\bullet_{U_\theta}\bigl(F^\bullet,\mathbb{D}(U_\theta)\bigr)[\dim X]\\
&=R\Hom_{U_\theta}\bigl(F^\bullet,\mathbb{D}(U_\theta)\bigr)[\dim X]
=\mathbb{D}_{alg}(V^\bullet)[\dim X].
\end{align*}}
由于 $L\Delta_\lambda$ 是 $D^b_{fg}(U_\theta)$ 与 $D^b_{coh}(\D_\lambda)$ 之间的等价，得到下述结果。

\begin{lemma}\label{lem:03.4.3}
设 $\lambda\in\h^*$ 正则。则下述函子图表交换
\begin{center}
\begin{tikzcd}[column sep=5.0em,row sep=2.9em]
D^b_{coh}(\D_\lambda) \arrow[r,"\mathbb{D}"] \arrow[d,"R\Gamma"'] & D^b_{coh}(\D_{-\lambda}) \arrow[d,"R\Gamma"]\\
D^b_{fg}(U_\theta) \arrow[r,"{\mathbb{D}_{alg}[\dim X]}"'] & D^b_{fg}(U_{-\theta})
\end{tikzcd}
\end{center}
\end{lemma}

设 $\alpha$ 是单根。若 $\lambda$ 是 $\alpha$-反支配的，则由 3.23 有 $R\Gamma(V^\bullet)=R\Gamma(LI_{s_\alpha}(V^\bullet))$。于是
{\small
\begin{align*}
R\Gamma(\mathbb{D}(V^\bullet))
&=\mathbb{D}_{alg}\bigl(R\Gamma(V^\bullet)\bigr)[\dim X]
=\mathbb{D}_{alg}\bigl(R\Gamma(LI_{s_\alpha}(V^\bullet))\bigr)[\dim X]\\
&=R\Gamma\bigl(\mathbb{D}(LI_{s_\alpha}(V^\bullet))\bigr).
\end{align*}}
这里 $\mathbb{D}(V^\bullet)\in D^b_{coh}(\D_{-\lambda})$，而 $\mathbb{D}(LI_{s_\alpha}(V^\bullet))\in D^b_{coh}(\D_{-s_\alpha\lambda})$。因此 $-s_\alpha\lambda$ 是 $\alpha$-反支配的。由此可见
$$R\Gamma(\mathbb{D}(V^\bullet))=R\Gamma\bigl(LI_{s_\alpha}\mathbb{D}(LI_{s_\alpha}(V^\bullet))\bigr).$$
由于 $\mathbb{D}(V^\bullet)$ 与 $LI_{s_\alpha}\mathbb{D}(LI_{s_\alpha}(V^\bullet))$ 都在 $D^b_{coh}(\D_{-\lambda})$ 中，且 $R\Gamma$ 是范畴等价，所以
$$\mathbb{D}(V^\bullet)=LI_{s_\alpha}\bigl(\mathbb{D}(LI_{s_\alpha}(V^\bullet))\bigr).$$
因此，在 $D^b_{coh}(\D_{-\lambda})$ 上成立
$$LI_{s_\alpha}\circ(\mathbb{D}\circ LI_{s_\alpha}\circ\mathbb{D})\cong 1.$$
这些函子都与几何平移交换，可见这个关系一般成立。于是对任意 $w\in W$ 有
$$LI_w\circ(\mathbb{D}\circ LI_{w^{-1}}\circ\mathbb{D})\cong 1.$$
这就证明了下述结果。

\begin{theorem}\label{thm:03.4.4}
交织函子 $LI_w:D^b_{coh}(\D_\lambda)\to D^b_{coh}(\D_{w\lambda})$ 的拟逆等于
$$\mathbb{D}\circ LI_{w^{-1}}\circ\mathbb{D}:D^b_{coh}(\D_{w\lambda})\to D^b_{coh}(\D_\lambda).$$
\end{theorem}

\section{不可约 $D_\lambda$-模的整体截面}\label{sec:03.3.5}

设 $\lambda\in\h^*$ 反支配。则对任意拟凝聚 $\D_\lambda$-模 $\mathcal{V}$，高次上同调模 $H^i(X,\mathcal{V})$（$i>0$）为零。因此只需研究 $\mathcal{V}$ 的整体截面 $\Gamma(X,\mathcal{V})$ 的行为。先给一个简单结果。

\begin{proposition}\label{prop:03.5.1}
设 $\lambda\in\h^*$ 反支配，$\mathcal{V}\in\mathcal{M}_{qc}(\D_\lambda)$ 不可约。则或者
\begin{enumerate}[label=(\roman*)]
\item $\Gamma(X,\mathcal{V})=0$；或者
\item $\mathcal{V}$ 由其整体截面 $\Gamma(X,\mathcal{V})$ 生成，且后者是不可约 $U_\theta$-模。
\end{enumerate}
\end{proposition}

\begin{proof}
如前所述，存在从 $\Delta_\lambda(\Gamma(X,\mathcal{V}))$ 映入 $\mathcal{V}$ 的自然态射，其像为凝聚 $\D_\lambda$-模。因此这个像等于 $0$ 或 $\mathcal{V}$。第一种情形 $\Gamma(X,\mathcal{V})=0$，(i) 成立。第二种情形 $\mathcal{V}$ 由其整体截面生成。还需证明 $\Gamma(X,\mathcal{V})$ 不可约。设 $K$ 是满态射 $\Delta_\lambda(\Gamma(X,\mathcal{V}))\to\mathcal{V}$ 的核。则有正合列
$$0\longrightarrow K\longrightarrow\Delta_\lambda(\Gamma(X,\mathcal{V}))\longrightarrow\mathcal{V}\longrightarrow 0,$$
从而
$$0\longrightarrow\Gamma(X,K)\longrightarrow\Gamma(X,\Delta_\lambda(\Gamma(X,\mathcal{V})))\longrightarrow\Gamma(X,\mathcal{V})\longrightarrow 0.$$
由 1.1，这说明 $\Gamma(X,K)=0$。

设 $U$ 是 $\Gamma(X,\mathcal{V})$ 的非零子模。则包含 $i:U\to\Gamma(X,\mathcal{V})$ 诱导同态 $\Delta_\lambda(i):\Delta_\lambda(U)\to\Delta_\lambda(\Gamma(X,\mathcal{V}))$。设 $\im\Delta_\lambda(i)$ 含于 $K$。施用 $\Gamma$ 便得到 $\Gamma(\Delta_\lambda(i))=0$，与 1.1 矛盾。因此 $\im\Delta_\lambda(i)$ 不含于 $K$。这说明自然态射 $K\oplus\Delta_\lambda(U)\to\Delta_\lambda(\Gamma(X,\mathcal{V}))$ 是满射。由 $\Gamma$ 的正合性可见，自然态射
{\small
\begin{align*}
\Gamma(X,K\oplus\Delta_\lambda(U))&=\Gamma(X,\Delta_\lambda(U))=U\\
&\longrightarrow\Gamma(X,\Delta_\lambda(\Gamma(X,\mathcal{V})))=\Gamma(X,\mathcal{V})
\end{align*}}
是满射；故 $U=\Gamma(X,\mathcal{V})$。
\end{proof}

上一结果有下述逆命题。

\begin{proposition}\label{prop:03.5.2}
设 $V$ 是 $\mathcal{M}(U_\theta)$ 中的不可约模\index{bukeyuemo@不可约模|hyperpage}，$\lambda\in\h^*$ 反支配。则存在不可约 $\D_\lambda$-模 $\mathcal{V}$，使 $\Gamma(X,\mathcal{V})\cong V$。这样的 $\D_\lambda$-模 $\mathcal{V}$ 在相差同构的意义下唯一。
\end{proposition}

证明从下述引理开始。

\begin{lemma}\label{lem:03.5.3}
设 $\lambda\in\h^*$ 反支配，$\mathcal{V}\in\mathcal{M}_{qc}(\D_\lambda)$。则 $\mathcal{V}$ 中存在不含非零整体截面的最大拟凝聚 $\D_\lambda$-子模 $\mathcal{V}'$。
\end{lemma}

\begin{proof}
设 $\mathcal{S}$ 是 $\mathcal{V}$ 的全体满足 $\Gamma(X,\mathcal{U})=0$ 的拟凝聚 $\D_\lambda$-子模所成族，并按包含关系把它看作有序集。设 $\mathcal{C}$ 是 $\mathcal{S}$ 中的链。由 [7, Ch.~I, 2.2.2]，$\mathcal{C}$ 中诸元的并 $\mathcal{W}$ 是 $\mathcal{V}$ 的拟凝聚 $\D_\lambda$-子模，且 $\Gamma(X,\mathcal{W})=0$。于是由 Zorn 引理，$\mathcal{S}$ 有极大元 $\mathcal{V}'$。设 $\mathcal{U}$ 是 $\mathcal{S}$ 的任一其他元。则 $\mathcal{U}+\mathcal{V}'$ 是 $\mathcal{V}$ 的拟凝聚 $\D_\lambda$-子模，且是 $\D_\lambda$-模 $\mathcal{U}\oplus\mathcal{V}'$ 的商。由 $\Gamma$ 的正合性可见 $\Gamma(X,\mathcal{U}+\mathcal{V}')=0$，即 $\mathcal{U}+\mathcal{V}'\in\mathcal{S}$。因此 $\mathcal{U}$ 是 $\mathcal{V}'$ 的 $\D_\lambda$-子模。
\end{proof}

现在可以证明 5.2。$V$ 的局部化 $\Delta_\lambda(V)$ 是凝聚 $\D_\lambda$-模，且由其整体截面生成。设 $\mathcal{W}$ 是 $\Delta_\lambda(V)$ 的、不等于 $\Delta_\lambda(V)$ 的凝聚 $\D_\lambda$-子模。则有凝聚 $\D_\lambda$-模的正合列
$$0\longrightarrow\mathcal{W}\longrightarrow\Delta_\lambda(V)\longrightarrow\mathcal{V}\longrightarrow 0,$$
且因 $\lambda$ 反支配，有
$$0\longrightarrow\Gamma(X,\mathcal{W})\longrightarrow\Gamma(X,\Delta_\lambda(V))\longrightarrow\Gamma(X,\mathcal{V})\longrightarrow 0.$$
由 1.1，$\Gamma(X,\Delta_\lambda(V))=V$；由假设 $V$ 不可约。于是 $\Gamma(X,\mathcal{W})$ 或为 $0$，或等于 $\Gamma(X,\Delta_\lambda(V))$。在第二种情形，$\Delta_\lambda(V)$ 的所有整体截面都已含于 $\mathcal{W}$。按 $\Delta_\lambda(V)$ 的定义，它作为 $\D_\lambda$-模由其整体截面生成。因此这会推出 $\mathcal{W}=\Delta_\lambda(V)$，与假设矛盾。故 $\Gamma(X,\mathcal{W})=0$，且 $\Gamma(X,\mathcal{V})=V$。于是由 5.3 可见，$\Delta_\lambda(V)$ 有最大凝聚 $\D_\lambda$-子模 $\mathcal{W}$，相应的 $\mathcal{V}$ 不可约；由 5.1，它由其整体截面生成。

最后证唯一性。设 $\mathcal{U}$ 是不可约 $\D_\lambda$-模，使 $\Gamma(X,\mathcal{U})=V$。则自然同态 $\Delta_\lambda(V)\to\mathcal{U}$ 的像或为 $0$，或为 $\mathcal{U}$。第一种情形 $\Gamma(X,\mathcal{U})=0$，与假设矛盾。因此这个同态是满射；由证明的第一部分，其核为 $\mathcal{W}$。5.2 证完。

下面要研究反支配的 $\lambda\in\h^*$ 时不可约 $\D_\lambda$-模的整体截面消没的充要条件。先讨论附着于 $\Pi_\lambda$ 中诸根所定反射的交织函子在不可约模上的作用。设 $\alpha\in\Pi_\lambda$。则
$$\Sigma^+_{s_\alpha}\cap\Sigma_\lambda=\Sigma^+_\lambda\cap\bigl(-s_\alpha(\Sigma^+_\lambda)\bigr)=\Sigma^+_\lambda\cap\bigl(-(\Sigma^+_\lambda-\{\alpha\})\cup\{\alpha\}\bigr)=\{\alpha\},$$
故由 3.21 知，$I_{s_\alpha}$ 的左上同调维数不超过 $1$。再设 $\lambda$ 反支配。则 $n(s_\alpha\lambda)\le 1$，由 C.3.1 可见 $\Gamma$ 在 $\mathcal{M}_{qc}(\D_{s_\alpha\lambda})$ 上的右上同调维数不超过 $1$。由 3.23，对任意 $\mathcal{V}\in\mathcal{M}_{qc}(\D_\lambda)$，有以 $H^p(X,L^qI_{s_\alpha}(\mathcal{V}))$ 为 $E_2$ 项、收敛到 $\Gamma(X,\mathcal{V})$ 的谱序列。于是这个谱序列在 $E_2$ 项处退化，且
\begin{enumerate}[label=(\alph*)]
\item $\Gamma(X,L_{-1}I_{s_\alpha}(\mathcal{V}))=H^1(X,I_{s_\alpha}(\mathcal{V}))=0$，
\item $U_\theta$-模 $\Gamma(X,\mathcal{V})$ 是 $\Gamma(X,I_{s_\alpha}(\mathcal{V}))$ 与 $H^1(X,L_{-1}I_{s_\alpha}(\mathcal{V}))$ 的扩张。
\end{enumerate}

\begin{lemma}\label{lem:03.5.4}
设 $\alpha\in\Pi_\lambda$，$\lambda\in\h^*$ 反支配，且 $\alpha$ 是 $\Pi_\lambda$ 中唯一与 $\lambda$ 正交的根。设 $\mathcal{V}\in\mathcal{M}_{qc}(\D_\lambda)$ 满足 $\Gamma(X,\mathcal{V})=0$。则
\begin{enumerate}[label=(\roman*)]
\item $I_{s_\alpha}(\mathcal{V})=0$；
\item $L_{-1}I_{s_\alpha}(\mathcal{V})=\mathcal{V}$。
\end{enumerate}
\end{lemma}

\begin{proof}
设 $\mu\in P(\Sigma)$ 是正则反支配权，$F$ 是以 $\mu$ 为最低权的不可约有限维 $\g$-模。记 $\mathcal{F}=\calO_X\otimes_\C F$。则如 C.2 所述，$U^\circ$-模 $G=(\mathcal{V}\otimes_{\calO_X}\mathcal{F})[\lambda+\mu]$ 带有自然滤过，其相应的分次模是 $\mathcal{V}(\nu)$ 之和，其中 $\nu$ 取遍 $F$ 的一切满足 $\lambda+\mu=w(\lambda+\nu)$（某 $w\in W$）的权。这个条件推出 $\lambda-w\lambda=w\nu-\mu\in Q(\Sigma)$，故 $w\in W_\lambda$。此外，因 $\lambda$ 反支配，此式左端是 $\Pi_\lambda$ 中诸根之和的相反数；又因 $\mu$ 是 $F$ 的最低权，右端是 $\Pi$ 中诸根之和。这只在两端都为 $0$ 时可能。因此 $w=1$ 或 $w=s_\alpha$。于是有正合列
$$0\longrightarrow\mathcal{V}(s_\alpha\mu)\longrightarrow G\longrightarrow\mathcal{V}(\mu)\longrightarrow 0.$$
此外，
$$H^i(X,G)=H^i(X,\mathcal{V}\otimes_{\calO_X}\mathcal{F})[\lambda+\mu]=\bigl(H^i(X,\mathcal{V})\otimes_\C F\bigr)[\lambda+\mu]=0,$$
由上同调的长正合列可见 $\Gamma(X,\mathcal{V}(s_\alpha\mu))=0$，$\Gamma(X,\mathcal{V}(\mu))=H^1(X,\mathcal{V}(s_\alpha\mu))$，且 $\mathcal{V}(s_\alpha\mu)$ 的高次上同调模消没。由 3.23，这最终说明
$$R\Gamma\bigl(LI_{s_\alpha}(\mathbb{D}(\mathcal{V}(\mu)))\bigr)=R\Gamma\bigl(\mathbb{D}(\mathcal{V}(s_\alpha\mu))[1]\bigr).$$
由导出范畴的等价（3.19）推得
$$LI_{s_\alpha}\bigl(\mathbb{D}(\mathcal{V}(\mu))\bigr)=\mathbb{D}(\mathcal{V}(s_\alpha\mu))[1],$$
即 $I_{s_\alpha}(\mathcal{V}(\mu))=0$，$L_{-1}I_{s_\alpha}(\mathcal{V}(\mu))=\mathcal{V}(s_\alpha\mu)$。(i) 与 (ii) 由 3.5(ii) 得出。
\end{proof}

\begin{lemma}\label{lem:03.5.5}
设 $\alpha\in\Pi_\lambda$，$\lambda\in\h^*$ 满足 $-\check{\alpha}(\lambda)=p\in\Z_+$。设 $\mathcal{V}$ 是拟凝聚 $\D_\lambda$-模，且 $I_{s_\alpha}(\mathcal{V})=0$，$L_{-1}I_{s_\alpha}(\mathcal{V})=\mathcal{V}(p\alpha)$。则 $\Gamma(X,\mathcal{V}(p\alpha))=0$。
\end{lemma}

\begin{proof}
由 3.23，对一切 $i\in\Z$ 有
$$H^i(X,\mathcal{V})=H^{i+1}\bigl(X,L_{-1}I_{s_\alpha}(\mathcal{V})\bigr)=H^{i+1}\bigl(X,\mathcal{V}(p\alpha)\bigr).$$
\end{proof}

\begin{lemma}\label{lem:03.5.6}
设 $\lambda\in\h^*$，$\alpha\in\Pi_\lambda$。置 $-\check{\alpha}(\lambda)=p\in\Z$。设 $\mathcal{V}$ 是不可约 $\D_\lambda$-模。则或者
\begin{enumerate}[label=(\roman*)]
\item $L_{-1}I_{s_\alpha}(\mathcal{V})=0$；或者
\item $I_{s_\alpha}(\mathcal{V})=0$，且 $L_{-1}I_{s_\alpha}(\mathcal{V})=\mathcal{V}(p\alpha)$。
\end{enumerate}
\end{lemma}

\begin{proof}
由 C.3.5(ii) 可设 $\lambda$ 反支配且正则。又因 $\mathcal{V}$ 不可约，$\Gamma(X,\mathcal{V})$ 是不可约 $U_\theta$-模，且或者 $\Gamma(X,I_{s_\alpha}(\mathcal{V}))=0$，或者 $H^1(X,L_{-1}I_{s_\alpha}(\mathcal{V}))=0$。由 C.4.3 推得：第一种情形 $I_{s_\alpha}(\mathcal{V})=0$，第二种情形 $L_{-1}I_{s_\alpha}(\mathcal{V})=0$。设 $I_{s_\alpha}(\mathcal{V})=0$。此时 $L_{-1}I_{s_\alpha}(\mathcal{V})\ne 0$。由 3.5(ii)，现在可设 $\lambda$ 反支配，且 $\alpha$ 是 $\Pi_\lambda$ 中唯一与其正交的根。此外，由前面的讨论与 $s_\alpha\lambda=\lambda$ 的反支配性推得
$$\Gamma(X,\mathcal{V})=H^1\bigl(X,L_{-1}I_{s_\alpha}(\mathcal{V})\bigr)=0.$$
断言由 3.5(ii) 与 4.4 得出。
\end{proof}

\begin{lemma}\label{lem:03.5.7}
设 $\lambda\in\h^*$，$\alpha\in\Pi_\lambda$。置 $-\check{\alpha}(\lambda)=p\in\Z$。设 $\mathcal{V}$ 是不可约 $\D_\lambda$-模，且 $L_{-1}I_{s_\alpha}(\mathcal{V})=0$。则 $I_{s_\alpha}(\mathcal{V})$ 有最大拟凝聚子模 $\mathcal{U}$，且
\begin{enumerate}[label=(\roman*)]
\item 有正合列
$$0\longrightarrow\mathcal{U}\longrightarrow I_{s_\alpha}(\mathcal{V})\longrightarrow\mathcal{V}(p\alpha)\longrightarrow 0;$$
\item $L_{-1}I_{s_\alpha}(\mathcal{U})=\mathcal{U}(-p\alpha)$，且当 $q\ne-1$ 时 $L_qI_{s_\alpha}(\mathcal{U})=0$。
\end{enumerate}
此外，若 $p\ge 0$，则 $\Gamma(X,\mathcal{U})=0$。
\end{lemma}

\begin{proof}
仍只需证第一个断言，此时可设 $\lambda$ 反支配且正则。如前所述，此时 $I_{s_\alpha}(\mathcal{V})$ 的一切高次上同调模消没。设 $\mathcal{W}$ 是 $I_{s_\alpha}(\mathcal{V})$ 的任一不等于 $I_{s_\alpha}(\mathcal{V})$ 的拟凝聚 $\D_{s_\alpha\lambda}$-子模。则由 C.3.1 与上同调的长正合列推得，$\mathcal{C}=I_{s_\alpha}(\mathcal{V})/\mathcal{W}$ 的高次上同调模也消没。因此由 C.4.3 推得 $\Gamma(X,\mathcal{C})\ne 0$。上同调的长正合列给出
$$0\longrightarrow\Gamma(X,\mathcal{W})\longrightarrow\Gamma(X,I_{s_\alpha}(\mathcal{V}))\longrightarrow\Gamma(X,\mathcal{C})\longrightarrow H^1(X,\mathcal{W})\longrightarrow 0.$$
\begin{annote}
此式中 $I_{s_\lambda}(\mathcal{V})$ 的下标 $s_\lambda$ 无意义（$\lambda$ 是权，不是根），应为 $I_{s_\alpha}(\mathcal{V})$；已按此译。[来源：底本 p.64，本证明通篇只涉及 $\alpha\in\Pi_\lambda$ 所定的反射 $s_\alpha$]
\end{annote}

可以取 $\nu\in P(\Sigma)$，使 $\lambda+\nu$ 反支配，且 $\alpha$ 是 $\Pi_\lambda$ 中唯一与 $\lambda+\nu$ 正交的根。于是 $\check{\alpha}(\lambda+\nu)=0$，且 $\check{\alpha}(\lambda)=-\check{\alpha}(\nu)=-p$。此时序列
$$0\longrightarrow\mathcal{W}(s_\alpha\nu)\longrightarrow I_{s_\alpha}(\mathcal{V})(s_\alpha\nu)\longrightarrow\mathcal{C}(s_\alpha\nu)\longrightarrow 0$$
正合；又因 $s_\alpha(\lambda+\nu)$ 反支配，得到正合列
$$0\longrightarrow\Gamma(X,\mathcal{W}(s_\alpha\nu))\longrightarrow\Gamma(X,I_{s_\alpha}(\mathcal{V})(s_\alpha\nu))\longrightarrow\Gamma(X,\mathcal{C}(s_\alpha\nu))\longrightarrow 0.$$
此外，由 3.5(ii) 与 3.23，
{\small
\begin{align*}
I_{s_\alpha}(\mathcal{V})(s_\alpha\nu)&=I_{s_\alpha}(\mathcal{V}(\nu)),\\
\Gamma(X,I_{s_\alpha}(\mathcal{V})(s_\alpha\nu))&=\Gamma(X,I_{s_\alpha}(\mathcal{V}(\nu)))=\Gamma(X,\mathcal{V}(\nu)),
\end{align*}}
故由 5.1，它或为 $0$，或是不可约 $U_\theta$-模。

现在断言 $\Gamma(X,\mathcal{C}(s_\alpha\nu))\ne 0$。设不然，即 $\Gamma(X,\mathcal{C}(s_\alpha\nu))=0$。则由 5.4 有
$$I_{s_\alpha}(\mathcal{C}(s_\alpha\nu))=0,\qquad L_{-1}I_{s_\alpha}(\mathcal{C}(s_\alpha\nu))=\mathcal{C}(s_\alpha\nu).$$
由 3.5(ii)，这推出
$$I_{s_\alpha}(\mathcal{C}(-p\alpha))=0,\qquad L_{-1}I_{s_\alpha}(\mathcal{C}(-p\alpha))=\mathcal{C}.$$
于是由 5.5，$\mathcal{C}$ 没有整体截面，与前面的讨论矛盾。这又推出 $\Gamma(X,\mathcal{W}(s_\alpha\nu))=0$；且由 5.4 与 5.5 可见 $\Gamma(X,\mathcal{W})=0$。于是已经证明：$I_{s_\alpha}(\mathcal{V})$ 的任一不等于 $I_{s_\alpha}(\mathcal{V})$ 的拟凝聚 $\D_{s_\alpha\lambda}$-子模 $\mathcal{W}$ 都满足 $\Gamma(X,\mathcal{W}(s_\alpha\nu))=0$。故由 5.3 可见，$I_{s_\alpha}(\mathcal{V})$ 含最大拟凝聚子模 $\mathcal{U}$，且 $\Gamma(X,\mathcal{U})=0$。此外，$(I_{s_\alpha}(\mathcal{V})/\mathcal{U})(s_\alpha\nu)$ 是不可约 $\D_{\lambda+\nu}$-模，且
$$\Gamma\bigl(X,(I_{s_\alpha}(\mathcal{V})/\mathcal{U})(s_\alpha\nu)\bigr)=\Gamma(X,\mathcal{V}(\nu)).$$
由 5.2 推得 $(I_{s_\alpha}(\mathcal{V})/\mathcal{U})(s_\alpha\nu)=\mathcal{V}(\nu)$。这意味着 $I_{s_\alpha}(\mathcal{V})/\mathcal{U}=\mathcal{V}(p\alpha)$。这就证明了 (i)。

前面已经得出 $\Gamma(X,\mathcal{U}(s_\alpha\nu))=0$。于是由 5.4，$L_{-1}I_{s_\alpha}(\mathcal{U}(s_\alpha\nu))=\mathcal{U}(s_\alpha\nu)$，且当 $q\ne-1$ 时 $L_qI_{s_\alpha}(\mathcal{U}(s_\alpha\nu))=0$。因此
$$L_{-1}I_{s_\alpha}(\mathcal{U})=\mathcal{U}(s_\alpha\nu-\nu)=\mathcal{U}(-p\alpha),\qquad
L_qI_{s_\alpha}(\mathcal{U})=0\quad(q\ne-1).$$
这就证明了 (ii)。由 (ii)，$L_{-1}I_{s_\alpha}(\mathcal{U}(-p\alpha))=\mathcal{U}$，其余导出交织函子在 $\mathcal{U}(-p\alpha)$ 上消没。由 5.5 可见，若 $p\ge 0$，则 $\Gamma(X,\mathcal{U})=0$。
\end{proof}

\begin{corollary}\label{cor:03.5.8}
设 $\lambda\in\h^*$，$\alpha\in\Pi_\lambda$，且 $-\check{\alpha}(\lambda)=p\in\Z_+$。设 $\mathcal{V}$ 是不可约 $\D_\lambda$-模，满足 $\Gamma(X,\mathcal{V})\ne 0$ 与 $L_{-1}I_{s_\alpha}(\mathcal{V})=0$。则 $\Gamma(X,\mathcal{V}(p\alpha))\ne 0$。
\end{corollary}

\begin{proof}
由 3.23，$\Gamma(X,I_{s_\alpha}(\mathcal{V}))=\Gamma(X,\mathcal{V})\ne 0$。于是由 $\Gamma$ 的左正合性与 5.7，$\Gamma(X,\mathcal{V}(p\alpha))\ne 0$。
\end{proof}

\begin{theorem}\label{thm:03.5.9}
设 $\lambda\in\h^*$ 反支配，$S$ 是 $\Pi_\lambda$ 中由与 $\lambda$ 正交的根组成的子集。设 $\mathcal{V}$ 是不可约 $\D_\lambda$-模。则下述条件等价：
\begin{enumerate}[label=(\roman*)]
\item $\Gamma(X,\mathcal{V})=0$；
\item 存在 $\alpha\in S$，使 $I_{s_\alpha}(\mathcal{V})=0$。
\end{enumerate}
\end{theorem}

\begin{proof}
(ii)$\Rightarrow$(i) 由 5.4，$L_{-1}I_{s_\alpha}(\mathcal{V})=\mathcal{V}$。于是由 5.5，$\Gamma(X,\mathcal{V})=0$。

(i)$\Rightarrow$(ii) 设 $W(\lambda)$ 是 $\lambda$ 在 $W$ 中的稳定化子。则 $W(\lambda)$ 由关于 $S$ 中诸根的反射生成。设 $F$ 是有限维表示，其最低权 $\nu$ 正则，并置 $\mathcal{F}=\calO_X\otimes_\C F$。则 $\mathcal{V}\otimes_{\calO_X}\mathcal{F}$ 满足
$$\Gamma(X,\mathcal{V}\otimes_{\calO_X}\mathcal{F})=\Gamma(X,\mathcal{V})\otimes_\C F=0.$$
此外 $\Gamma\bigl(X,(\mathcal{V}\otimes_{\calO_X}\mathcal{F})[\lambda+\nu]\bigr)=0$。另一方面，考虑 C.2 中讨论的 $F$ 的滤过，它诱导 $(\mathcal{V}\otimes_{\calO_X}\mathcal{F})[\lambda+\nu]$ 的滤过，其相应的分次层是 $\mathcal{V}(\mu)$ 的直和，其中 $\mu$ 取遍 $F$ 的一切满足 $w(\lambda+\nu)=\lambda+\mu$（某 $w\in W$）的权。这说明 $w\lambda-\lambda=\mu-w\nu$，故 $w\in W_\lambda$。此外，等式 $\lambda-w^{-1}\lambda=w^{-1}\mu-\nu$ 的左端是 $\Pi_\lambda$ 中诸根之和的相反数，右端是 $\Pi$ 中诸根之和。这说明 $w\lambda=\lambda$，即 $w\in W(\lambda)$。设 $w_1\in W(\lambda)$ 使 $\mathcal{V}(w_1\nu)$ 是 $(\mathcal{V}\otimes_{\calO_X}\mathcal{F})[\lambda+\nu]$ 的 $\calO_X$-子模，则 $\Gamma(X,\mathcal{V}(w_1\nu))=0$。现在设对一切 $\alpha\in S$ 都有 $L_{-1}I_{s_\alpha}(\mathcal{V})=0$。断言此时对任意 $w\in W(\lambda)$ 有 $\Gamma(X,\mathcal{V}(w\nu))\ne 0$，与假设矛盾。对 $w$ 在 $W(\lambda)$ 中的长度作归纳；由 [5, Ch.~IV, \S1, no.~8, Cor.~4 of Prop.~7]，这个长度等于 $w$ 在 $W_\lambda$ 中的长度。若 $w=1$，$\lambda+\nu$ 反支配且正则，结论由 C.4.4 得出。设 $\ell_\lambda(w)=k>0$。取 $w=s_\alpha w'$，其中 $\alpha\in S$，$w'\in W(\lambda)$，$\ell_\lambda(w')=k-1$。则由 [5, Ch.~VI, \S1, no.~6, Cor.~1 of Prop.~17] 可见 $w'^{-1}\alpha\in\Sigma^+_\lambda$。由 $\nu$ 的反支配性，这推出
$$\check{\alpha}(\lambda+w'\nu)=\check{\alpha}(w'\nu)=\bigl(w'^{-1}\alpha\bigr)^{\vee}(\nu)\in-\N,$$
且 $\lambda+w'\nu$ 是 $\alpha$-反支配的。由归纳假设，$\Gamma(X,\mathcal{V}(w'\nu))\ne 0$；由 3.5(ii) 又有
$$L_{-1}I_{s_\alpha}(\mathcal{V}(w'\nu))=L_{-1}I_{s_\alpha}(\mathcal{V})(w\nu)=0.$$
因此断言由 5.8 得出。
\end{proof}

\section{交织函子与完整\index{wanzheng@完整|hyperpage}复形\index{wanzhengfuxing@完整复形|hyperpage}}\label{sec:03.3.6}

完整模所成范畴是范畴 $\mathcal{M}_{qc}(\D_\lambda)$ 的厚子范畴。因此可以考虑范畴 $D^b_{hol}(\D_\lambda)$，即 $D^b(\D_\lambda)$ 中由上同调完整的复形所成的满子范畴。显然，几何平移函子 $V\mapsto V(\mu)$（$\mu\in P(\Sigma)$）诱导范畴 $D^b_{hol}(\D_\lambda)$ 与范畴 $D^b_{hol}(\D_{\lambda+\mu})$ 的等价。

每个完整模都有有限长。由此得到下述结果。

\begin{lemma}\label{lem:03.6.1}
不可约完整 $\D_\lambda$-模在 $D^b_{hol}(\D_\lambda)$ 中构成生成类。
\end{lemma}

\begin{proof}
按标准的截断论证 [1, I.12.6]，全体完整 $\D_\lambda$-模在 $D^b_{hol}(\D_\lambda)$ 中构成生成类。另一方面，设 $\mathcal{V}$ 是长为 $n$ 的完整 $\D_\lambda$-模，$\mathcal{U}$ 是它的一个极大凝聚子模。则有正合列
$$0\longrightarrow\mathcal{U}\xrightarrow{\ i\ }\mathcal{V}\longrightarrow\mathcal{V}/\mathcal{U}\longrightarrow 0,$$
其中 $\mathcal{V}/\mathcal{U}$ 不可约，$\mathcal{U}$ 的长为 $n-1$。设下述特异三角附着于态射 $i$
\begin{center}
\begin{tikzcd}[column sep=3.4em,row sep=2.2em]
C_i \arrow[rr] \arrow[dr] & & \mathbb{D}(\mathcal{U}) \arrow[dl,"\mathbb{D}(i)"']\\
& \mathbb{D}(\mathcal{V}) \arrow[ul,"{[1]}"'] &
\end{tikzcd}
\end{center}
则锥 $C_i$ 同构于 $\mathbb{D}(\mathcal{V}/\mathcal{U})$。因此对长度作归纳可见，由不可约完整 $\D_\lambda$-模生成的 $D^b_{hol}(\D_\lambda)$ 的三角子范畴包含一切形如 $\mathbb{D}(\mathcal{V})$ 的复形，其中 $\mathcal{V}$ 是完整 $\D_\lambda$-模。由第一段注记，这个子范畴等于 $D^b_{hol}(\D_\lambda)$。
\end{proof}

按定义，对任意 $w\in W$，交织函子 $LI_w$ 把 $D^b_{hol}(\D_\lambda)$ 映入 $D^b_{hol}(\D_{w\lambda})$。

\begin{theorem}\label{thm:03.6.2}
设 $w\in W$，$\lambda\in\h^*$。则 $LI_w$ 是范畴 $D^b_{hol}(\D_\lambda)$ 与 $D^b_{hol}(\D_{w\lambda})$ 的等价。
\end{theorem}

\begin{proof}
由乘积公式（3.18），只需对单反射证明本命题。利用几何平移与 3.5，还可设 $\lambda$ 反支配且正则。此时有两种可能：
\begin{enumerate}[label=(\alph*)]
\item $\check{\alpha}(\lambda)\notin-\N$；
\item $\check{\alpha}(\lambda)\in-\N$。
\end{enumerate}
在情形 (a)，$s_\alpha\lambda$ 仍正则反支配。因此由 3.22，$I_{s_\alpha}$ 是正合函子，且以 $I_{s_\alpha^{-1}}$ 为逆。这立即给出断言。

在情形 (b)，$p=-\check{\alpha}(\lambda)\in\N$。此时只需证明：$D^b_{hol}(\D_{s_\alpha\lambda})$ 中由同构于复形 $LI_{s_\alpha}(C^\bullet)$（$C^\bullet\in D^b_{hol}(\D_\lambda)$）的对象所成的满子范畴 $\mathcal{A}$ 等于 $D^b_{hol}(\D_{s_\alpha\lambda})$。为此，由 6.1，只需证明 $\mathcal{A}$ 含一切复形 $\mathbb{D}(\mathcal{V}(p\alpha))$，其中 $\mathcal{V}$ 取遍不可约完整 $\D_\lambda$-模。

由 5.6，对不可约完整 $\D_\lambda$-模 $\mathcal{V}$ 有两种可能：
\begin{enumerate}[label=(\roman*)]
\item $L_{-1}I_{s_\alpha}(\mathcal{V})=\mathcal{V}(p\alpha)$，且 $I_{s_\alpha}(\mathcal{V})=0$；
\item $L_{-1}I_{s_\alpha}(\mathcal{V})=0$。
\end{enumerate}
\begin{annote}
底本在 §6 的证明中把 §5 的条目引作「4.1」「4.6」「4.7」「4.9」，但 §4 只有 4.1–4.4，且 4.1 是扭变与对偶交换的引理，与有限长无关。按内容，四处分别是 5.1、5.6、5.7、5.9；此处及 6.2、6.3 的证明均已按后者译。[来源：底本 p.66–67，与 5.1、5.6、5.7、5.9 的内容逐一对照]
\end{annote}
若 (i) 成立，则有 $LI_{s_\alpha}(\mathbb{D}(\mathcal{V}))=\mathbb{D}(\mathcal{V}(p\alpha))[1]$。因此 $LI_{s_\alpha}(\mathbb{D}(\mathcal{V})[-1])=\mathbb{D}(\mathcal{V}(p\alpha))$；于是对一切 (i) 型不可约 $\D_\lambda$-模 $\mathcal{V}$，复形 $\mathbb{D}(\mathcal{V}(p\alpha))$ 都在 $\mathcal{A}$ 中。这也说明，若完整 $\D_\lambda$-模 $U$ 的合成因子全为 (i) 型，则复形 $\mathbb{D}(U(p\alpha))$ 也在 $\mathcal{A}$ 中。

若 (ii) 成立，则由 5.7 有短正合列
$$0\longrightarrow\mathcal{U}\longrightarrow I_{s_\alpha}(\mathcal{V})\longrightarrow\mathcal{V}(p\alpha)\longrightarrow 0.$$
可以取 $\nu\in P(\Sigma)$，使 $\lambda+\nu$ 反支配，且 $\alpha$ 是 $\Pi_\lambda$ 中唯一与 $\lambda+\nu$ 正交的根。于是由 3.5，
$$0\longrightarrow\mathcal{U}(s_\alpha\nu)\longrightarrow I_{s_\alpha}(\mathcal{V}(\nu))\longrightarrow\mathcal{V}(\nu)\longrightarrow 0$$
正合。此外由 5.7 可见 $\Gamma(X,\mathcal{U}(s_\alpha\nu))=0$。因 $\lambda+\nu$ 反支配，由 C.3.2，$\Gamma$ 正合，故 $\mathcal{U}(s_\alpha\nu)$ 的一切合成因子都没有整体截面。由 5.9，这说明 $\mathcal{U}(s_\alpha\nu)$ 的一切合成因子都是 (i) 型。由 3.5 推得，$\mathcal{U}$ 的一切合成因子都是 (i) 型。由证明的第一部分可见 $U$ 在 $\mathcal{A}$ 中。现在考虑附着于态射 $I_{s_\alpha}(\mathcal{V})\to\mathcal{V}(p\alpha)$ 的特异三角
\begin{center}
\begin{tikzcd}[column sep=3.8em,row sep=2.4em]
C^\bullet \arrow[rr] \arrow[dr] & & LI_{s_\alpha}(\mathbb{D}(\mathcal{V})) \arrow[dl]\\
& \mathbb{D}(\mathcal{V}(p\alpha)) \arrow[ul,"{[1]}"'] &
\end{tikzcd}
\end{center}
锥 $C^\bullet$ 同构于 $\mathbb{D}(U)[-1]$。由于 $LI_{s_\alpha}(\mathbb{D}(\mathcal{V}))$ 与 $\mathbb{D}(U)[-1]$ 都在 $\mathcal{A}$ 中，可见 $\mathbb{D}(\mathcal{V}(p\alpha))$ 也在 $\mathcal{A}$ 中。由 6.1 得 $\mathcal{A}=D^b_{hol}(\D_{s_\alpha\lambda})$。
\end{proof}

这个结果有下述推论。

\begin{theorem}\label{thm:03.6.3}
设 $\lambda\in\h^*$，$\theta=W\cdot\lambda$。设 $\mathcal{V}$ 是完整 $\D_\lambda$-模。则 $H^p(X,\mathcal{V})$（$p\in\Z_+$）是有限长 $U_\theta$-模。
\end{theorem}

\begin{proof}
取 $\mu\in\theta$ 反支配，$w\in W$ 使 $\lambda=w\mu$。由 6.2，存在上同调完整的复形 $C^\bullet$，使 $LI_w(C^\bullet)=\mathbb{D}(\mathcal{V})$。由于 $\Gamma$ 是从 $\mathcal{M}_{qc}(\D_\mu)$ 映入 $\mathcal{M}(U_\theta)$ 的正合函子，可见 $R\Gamma(C^\bullet)=\Gamma(C^\bullet)$，且对任意 $p\in\Z$ 有 $H^p(R\Gamma(C^\bullet))=\Gamma(X,H^p(C^\bullet))$。由 5.1 还可推出，$H^p(R\Gamma(C^\bullet))$（$p\in\Z$）是有限长 $U_\theta$-模。因此 $R\Gamma(C^\bullet)$ 是上同调模有限长的 $U_\theta$-模复形。最后由 3.23 得到
$$H^p(X,\mathcal{V})=H^p\bigl(R\Gamma(\mathbb{D}(\mathcal{V}))\bigr)=H^p\bigl(R\Gamma(LI_w(C^\bullet))\bigr)=H^p\bigl(R\Gamma(C^\bullet)\bigr).$$
\end{proof}

\begin{proposition}\label{prop:03.6.4}
设 $\lambda\in\h^*$ 正则反支配，$\theta=W\cdot\lambda$。设 $V$ 是有限生成 $U_\theta$-模。则下述条件等价：
\begin{enumerate}[label=(\roman*)]
\item $\Delta_\lambda(V)$ 是完整 $\D_\lambda$-模；
\item 存在 $w\in W$，使 $L\Delta_{w\lambda}(V)$ 是上同调完整的复形；
\item $H^p(\n_x,V)$（$p\in\Z_+$）对一切 $x\in X$ 有限维；
\item 存在 $w\in W$，使 $H^p(\n_x,V)(w\lambda+\rho)$（$p\in\Z_+$）对一切 $x\in X$ 有限维。
\end{enumerate}
\end{proposition}

\begin{proof}
先注意，由 1.19，对任意 $w\in W$ 有 $L\Delta_{w\lambda}(V)\in D^b_{coh}(\D_{w\lambda})$。由 3.19 与 6.2，(i) 与 (ii) 等价。又 (iii) 蕴含 (iv)。设 (ii) 成立。对 $x\in X$，记 $i_x$ 为 $x$ 到 $X$ 的包含。则
$$LT_x\bigl(L\Delta_{w\lambda}(V)\bigr)=Li_x^*\bigl(L\Delta_{w\lambda}(V)\bigr)$$
是上同调有限维的复形。由 2.6，这说明 (iv)。反之，若 (iv) 成立，则同一结果说明复形 $LT_x(L\Delta_{w\lambda}(V))$ 对一切 $x\in X$ 都有有限维上同调。于是由（[1]，VII.10.7.(ii)）推得 (ii)。
\end{proof}

\begin{corollary}\label{cor:03.6.5}
设 $\theta$ 是 $\h^*$ 中由正则元素构成的 $W$-轨道。设 $V$ 是有限生成 $U_\theta$-模。若 $H^p(\n_x,V)$（$p\in\Z_+$）对一切 $x\in X$ 有限维，则模 $V$ 有有限长。
\end{corollary}

\begin{proof}
取 $\lambda\in\theta$ 反支配。则由 6.4 可见 $\Delta_\lambda(V)$ 完整，因而有有限长。由范畴的等价，$V$ 有有限长。
\end{proof}
